% Pillar-1 (scaling) discussion + conclusion
\section{Discussion and Limitations}
\label{sec:p1-discussion}
The most useful result here is a boundary on what can be claimed. On this
benchmark, neither mean GRPO training reward $\bar{R}$ nor fitted endpoint ceiling $R_{\max}$ exhibits a measurable, monotone dependence
on model size, and no universal saturation exponent is identifiable; the honest
description is a local, stack-conditioned taxonomy of endpoint behavior rather than a
Chinchilla-style power law in $N$. Several methodological limitations qualify even this modest
statement:
\begin{enumerate}[leftmargin=*,nosep]
  \item \textbf{Step-1 pre-saturation on GSM8K.} Modern instruction-tuned checkpoints (`Qwen3.5-4B`, `Llama-3.1-8B-Instruct`, `DeepSeek-V3.1`, `Kimi-K2-Thinking`, `Qwen3-30B-MoE-Inst`, `Qwen3-235B-MoE`) already achieve $R(1) \in [0.875, 1.00]$ on GSM8K at step~1 (\tableref{tab:p1-rawtrace}). Consequently, fitting $R(t) = R_{\max}(1 - e^{-\lambda t})$ with $R(0)=0$ forces $\lambda$ to the optimizer upper bound ($\lambda = 10$), and even with a free baseline offset $c \in [0,1]$ the remaining headroom $(1 - c)$ is smaller than the binomial rollout noise floor $\sigma_{\text{step}} \in [0.074, 0.119]$ with $8$--$16$ rollouts per step over $T \le 30$ steps. Evaluating post-training scaling laws on harder, unsaturated reasoning benchmarks (such as AIME, OlympiadBench, or competition coding) over $T > 200$ steps is essential before concluding whether GRPO rate constants scale with parameter count.
  \item \textbf{Single-seed frontier API traces and summary-reconstructed diagnostics.} The frontier-scale anchors are single-seed \textsc{Tinker} API runs and must be read as descriptive case studies rather than powered multi-seed replications. Earlier phase-score, functional-form, truncation, iso-compute and phase-conformity batteries ran on traces synthesised from six per-anchor summary statistics; they have been withdrawn and replaced by raw-trace recomputations (\secref{sec:p1-rawtrace}).  Five extended-pool probes (Qwen3-32B, Qwen3.5-27B, Qwen3-30B-MoE, Qwen3-30B-MoE-Inst, Qwen3-235B-MoE) are interrupted runs with 3--5 console-transcribed rewards; they are reported descriptively and excluded from fits.  Inferential statements are restricted to the seven full traces (20--30 steps) and to the 40-step same-stack runs, and even there the unit is a single-seed anchor.
  \item \textbf{Endogeneity of post-hoc capability clustering and variance starvation.} Partitioning anchors by Ward $k=2$ clustering on $R_{\max}$ (\S\ref{par:scaling-iter133}) is a descriptive variance decomposition of the outcome distribution rather than an exogenous causal predictor. The same partition aligns with the step-1 baseline $R(1) \ge 0.50$ (11 of 12 anchors) and with the zero-variance fraction (\textsc{zvf}), but both are measured on the same runs and the threshold was chosen after the fact, so neither is an independent covariate; an a-priori indicator (e.g.\ base vs.\ instruction-tuned checkpoint, or a zero-shot held-out evaluation before RL) on a larger pool is needed. Disentangling parameter scale from within-group contrast starvation requires active curriculum or rollout filtering, such as dynamic sampling in DAPO~\citep{yu2025dapo}, unbiased advantage normalization in Dr.\ GRPO~\citep{tong2025drgrpo}, or online difficulty filtering~\citep{bae2025online, jiang2025vcrl}.
\end{enumerate}

\paragraph{Submission scope.}
For a workshop-short / limits-audit venue we claim only the descriptive flat
slope, the absence of the clean three-phase template on the full traces, and
non-identifiability of a universal saturation exponent on the released
single-seed-heavy corpus. We do \emph{not} claim a powered multi-seed
cross-scale test; that remains future work and is not required for the
identifiability boundary already established here.

\section{Conclusion}
\label{sec:p1-conclusion}
We asked whether GRPO post-training obeys a clean scaling law in model size and
found, on a controlled cross-library and cross-scale benchmark, that it does not---at least not
one identifiable from $2.4$ orders of magnitude ($4$B--$1$T parameters) of short-horizon GSM8K evidence. The
cross-scale mean-reward and ceiling slopes are not distinguishable from zero (with intervals too wide to
establish flatness), no full trace follows a universal three-phase saturation template, and the Nemotron-3-Super-120B-A12B LatentMoE trace is best treated
as a format-collapse structural outlier ($55\%$ zero-reward steps). What survives is a taxonomic, endpoint-level
description of the ceiling, dominated by pre-existing instruction-following capability and rollout contrast rather
than an added $\log_{10} N$ term, and a concrete experimental roadmap:
re-plotting post-training scaling on usable contrastive compute $C_{\text{eff}}$, extending horizons beyond $T^\star > 200$ steps on unsaturated benchmarks, and running multi-seed matched-budget evaluations. The released traces, telemetry, adapters, and scripts make that follow-up directly reproducible.
